% !PRINTED 35
\weberpage{35}
%<<CONT
Grad ein Theiler von $a$ ist, ist eine in $S$ enthaltene Gruppe; denn sind $A, A'$ zwei solche Elemente, so ist nach 3. der Grad von $AA'$ ein Theiler von $a$, und $AA'$ ist also auch in $\mathfrak{A}$ enthalten.

Ebenso ist der Inbegriff $\mathfrak{B}$ der Elemente $B$, deren Grad ein Theiler von $b$ ist, eine Gruppe. Das Element 1 kommt sowohl in $\mathfrak{A}$ als in $\mathfrak{B}$ vor, sonst enthalten beide kein gemeinschaftliches Element.

Der Grad $a'$ von $\mathfrak{A}$ ist relativ prim zu $b$. Denn ist $r$ eine in $a'$ aufgehende Primzahl, so giebt es nach 2. in $\mathfrak{A}$ ein Element vom Grade $r$. Da aber der Grad jedes Elementes von $\mathfrak{A}$ ein Theiler von $a$ ist, so ist auch $r$ ein Theiler von $a$ und nicht von $b$.

Ebenso beweist man, dass der Grad $b'$ von $\mathfrak{B}$ relativ prim zu $a$ ist.

Nun bestimme man (nach Bd.\,I, \S.\,118) zwei ganze Zahlen $x$ und $y$, so dass
\weq{6}{a x + b y = 1}
ist, und nehme irgend ein Element $\Theta$ von $S$. Dann ist
\weq{7}{\Theta = \Theta^{ax} \Theta^{by},}
und nun ist, da $(\Theta^{by})^a = 1$ ist, $\Theta^{by}$ in $\mathfrak{A}$ und aus dem gleichen Grunde $\Theta^{ax}$ in $\mathfrak{B}$ enthalten. Also folgt aus (7):
\weq{8}{\Theta = AB.}
Demnach ist jedes Element $\Theta$ in der Form $AB$ enthalten. Eine solche Darstellung ist aber auch nur auf eine Art möglich. Denn sind $A', B'$ zwei Elemente aus $\mathfrak{A}$ und $\mathfrak{B}$, und ist
\[
AB = A'B',
\]
so folgt, wenn man in die Potenz $by = 1 - ax$ erhebt, $A = A'$ und folglich auch $B = B'$. Die Anzahl der verschiedenen Producte der Form $AB$ ist aber $= a'b'$, und daher hat man
\[
n = ab = a'b',
\]
und da $a$ relativ prim zu $b'$ und $b$ relativ prim zu $a'$ ist:
\[
a' = a,\quad b' = b.
\]
Damit ist der Satz 4. in allen seinen Theilen bewiesen.

Wenn nun die beiden Gruppen $\mathfrak{A}, \mathfrak{B}$ durch Basen dargestellt sind, so folgt aus der Formel (8), dass auch $S$ durch eine Basis dargestellt ist, und die Basis von $S$ enthält die Elemente der Basen von $\mathfrak{A}$ und $\mathfrak{B}$ und keine anderen.
